\documentclass[11pt]{article}

% --- Packages ---
\usepackage[left=0.5in, right=0.5in, top=1in, bottom=1in]{geometry}
\usepackage{booktabs}       % Professional table rules
\usepackage{amsmath}         % Math environments
\usepackage{amssymb}         % Additional math symbols
\usepackage[T1]{fontenc}
\usepackage{lmodern}         % Latin Modern font

% Remove page number for a standalone supplementary table
\pagestyle{empty}

\begin{document}

\begin{center}
{\large\bfseries Supplementary Table 1}
\end{center}

\medskip
\noindent
Paired comparison of local largest Lyapunov exponent (LLE) during stimulation vs.\ no-stimulation periods across adaptation conditions.

\bigskip

\begin{table}[h]
\centering
\begin{tabular}{l c r r r r r}
\toprule
Condition & Networks & Wilcoxon $p$ & Cohen's $d$ & Median $\Delta$LLE & Mean Stim LLE & Mean No-Stim LLE \\
\midrule
No Adaptation & 25 & 0.035\rlap{*}  &  0.41 &  0.46 &  0.52 & 0.25 \\
SFA Only       & 25 & 0.001\rlap{*}  &  0.80 &  0.38 &  0.41 & 0.14 \\
STD Only       & 25 & 0.313   & $-$0.21 & $-$0.08 & $-$0.13 & 0.00 \\
SFA + STD      & 25 & 0.476   & $-$0.25 &  0.04 & $-$0.02 & 0.03 \\
\bottomrule
\end{tabular}
\end{table}

\smallskip
\noindent
\textit{Note.}\ $N = 300$ neurons per network, connection density = 33\%, $f \in [0.40,\, 0.60]$.
Wilcoxon signed-rank test (two-tailed) compared mean local LLE during stimulation
versus no-stimulation periods.
$\Delta$LLE = LLE$_{\text{stim}}$ $-$ LLE$_{\text{no-stim}}$.\\[2pt]
* $p < .05$.

\end{document}
